\documentclass[11pt,french]{article}


\usepackage{amsmath,amsfonts,graphicx}
\renewcommand{\baselinestretch}{1.1}

\usepackage{enumerate,multirow}
\usepackage{verbatim}
%\usepackage{eclbip}
%\usepackage{pst-node}
%\include{psfig.sty}
\usepackage{url}
\usepackage{fancyhdr}
%\usepackage{slashbox}
\usepackage[colorlinks=true]{hyperref}

%\usepackage{soul}

\usepackage{ulem}

\parindent=0cm
\textwidth 6.25in
\textheight 21.9 cm
\topmargin -1.4 cm
\oddsidemargin 0in
\evensidemargin 0in

%%%% debut macro anti-orphelins %%%%
%\widowpenalty=10000
%\clubpenalty=10000
%\raggedbottom
%%%% fin macro %%%%

\headsep 1.5cm
\footskip 1.5cm
\fancyhead{}
\fancyfoot{}
\renewcommand{\headrulewidth}{0.6pt}
%\renewcommand{\footrulewidth}{0.6pt}
\fancyhead[LO,LE]{{\footnotesize Universit\'e Laval \\
Facult\'e des sciences et de g\'enie\\
D\'epartement de math\'ematiques et de statistique
}}
\fancyhead[RO,RE]{{\footnotesize Automne 2026\\STT-2902 Modélisation statistique\\Julien Miron}}
%\fancyfoot[LO, LE] {\includegraphics[width=1.5cm]{UL.jpg}}
%\fancyfoot[RO, RE] {\vspace{-0.5cm} \thepage}


\pagestyle{fancy}
\pagenumbering{gobble}


\begin{document}

\begin{center}
    \section*{Présentation en équipe}

\end{center}

\subsection*{Description}
Six présentations en équipe auront lieu durant la session d'automne 2026. Chaque équipe se voit assigner un \textbf{sujet} ainsi qu'un \textbf{style de présentation imposé}. Les consignes détaillées (sujet, style, requis spécifiques) seront remises à chaque équipe environ deux semaines avant sa présentation.

\begin{center}
\begin{tabular}{|c|l|l|}
    \hline
    \textbf{Équipe} & \textbf{Sujet} & \textbf{Style imposé} \\
    \hline
    1 & Statistiques descriptives & Jeu télévisé / quiz interactif \\
    2 & Estimation ponctuelle & Cours magistral \\
    3 & Intervalles de confiance & Atelier expérimental \\
    4 & Tests d'hypothèses & Procès simulé / débat \\
    5 & Régression linéaire & Séance de révision d'examen \\
    6 & Tests d'indépendance & Enquête en classe et analyse en direct \\
    \hline
\end{tabular}
\end{center}

Votre tâche consiste à :
\begin{itemize}
    \item \textbf{Préparer une présentation} dynamique et pédagogique du sujet assigné, \textbf{dans le style imposé} à votre équipe;
    \item \textbf{Expliquer clairement les concepts} importants liés au sujet;
    \item \textbf{Engager les autres étudiants} selon le format prévu par votre style de présentation;
    \item \textbf{Faire le lien avec des situations réelles} ou des applications concrètes;
    \item \textbf{Travailler en équipe} pour bien répartir les tâches et assurer une présentation fluide et cohérente.
\end{itemize}

\subsection*{Consignes}

\begin{itemize}
    \item Le style de présentation imposé doit être respecté; une présentation qui s'écarte du style assigné sera pénalisée;
    \item La présentation doit durer entre 30 et 50 minutes;
    \item Tous les membres doivent prendre part activement à la présentation, avec un temps de parole comparable;
    \item Tous les éléments \textbf{requis} de la fiche de votre équipe doivent être couverts;
    \item Un support visuel (diapos, tableau, matériel, etc.) est obligatoire et doit être adapté au style imposé;
    \item Le matériel de présentation (diapos, questions de quiz, données, etc.) doit être remis à l'enseignant au moins \textbf{48 heures avant} la présentation;
    \item Adoptez une posture d’enseignant envers vos collègues : assurez un bon équilibre entre une présentation vivante et un rythme qui permet à votre auditoire de prendre des notes au besoin.

\end{itemize}

\subsection*{Évaluation}

L’évaluation tiendra compte des critères suivants :
\begin{itemize}
    \item Respect du style de présentation imposé et des requis spécifiques;
    \item Qualité de l’explication et maîtrise du sujet;
    \item Clarté et organisation de la présentation;
    \item Qualité de l’animation;
    \item Collaboration entre les membres de l’équipe;
    \item Participation lors des activités des autres équipes;
    \item Contribution au travail d'équipe;
    \item Évaluation par les pairs.
\end{itemize}

\end{document}